\documentclass[11pt,a4paper]{article}
\usepackage[margin=25mm]{geometry}
\usepackage[T1]{fontenc}
\usepackage[utf8]{inputenc}
\usepackage{lmodern,microtype,amsmath,amssymb,amsthm,booktabs,longtable,enumitem}
\usepackage[round,sort]{natbib}
\usepackage{xurl}
\usepackage[colorlinks=true,allcolors=blue]{hyperref}
\setlength{\emergencystretch}{2em}
\numberwithin{equation}{section}
\newtheorem{theorem}{Theorem}[section]
\newtheorem{proposition}[theorem]{Proposition}
\newtheorem{lemma}[theorem]{Lemma}
\newtheorem{corollary}[theorem]{Corollary}
\newtheorem{conjecture}[theorem]{Conjecture}
\newtheorem{maintheorem}[theorem]{Theorem}
\theoremstyle{definition}
\newtheorem{definition}[theorem]{Definition}
\newtheorem{fact}{Standard fact}
\renewcommand{\thefact}{I\arabic{fact}}
\theoremstyle{remark}
\newtheorem{remark}[theorem]{Remark}
\newcommand{\Fp}{\mathbb F_p}
\newcommand{\Fq}{\mathbb F_{p^2}}
\newcommand{\Fb}{\overline{\mathbb F}_p}
\newcommand{\Z}{\mathbb Z}
\newcommand{\Q}{\mathbb Q}
\newcommand{\Res}{\operatorname{Res}}
\newcommand{\Aut}{\operatorname{Aut}}
\newcommand{\Stab}{\operatorname{Stab}}
\newcommand{\den}{\operatorname{den}_{\mathrm{red}}}
\newcommand{\leg}[2]{\left(\frac{#1}{#2}\right)}
\newcommand{\Lp}[1]{\left(\frac{#1}{p}\right)}
% Set this to the public evidence repository URL or DOI landing page before submission.
\newcommand{\DataRepositoryURL}{https://github.com/bbrhuft/factorial-supersingular-congruence/tree/main/supplementary}

% ---------------------------------------------------------------------
% v2.0 (9 Nov 2026). Paper A of two. The combined manuscript v1.14 was split:
% this paper contains the factorial evaluation (old Sections 1.2, 2, 3) and
% its computational evidence; the Fricke application is in the companion
% paper "Supersingular poles and signed residues of prime-level Fricke
% Hauptmoduln". Statements and proofs are those of v1.14, except:
%  * new Remark on eta coordinates in Section 3 (moved from old Section 1.3);
%  * citations to Robert (1980) checked against the paper and extended.
% Open items for the author: search for "TODO(author)".
% ---------------------------------------------------------------------
\title{A factorial supersingular congruence\\and signed discriminant ratios}
\author{David Jordan}
\date{9 November 2026}
\begin{document}
\maketitle
\begin{abstract}
For every prime $p\ge5$, we prove the supersingular congruence
\[
E_{p+1}^{(p-1)/2}\equiv
\left(\frac6p\right)\left(\frac{p-1}{2}\right)!
\Delta^{(p^2-1)/24}\pmod{(p,E_{p-1})}.
\]
It evaluates a half-power ratio uniformly on the supersingular locus
and recovers Baker's congruence upon squaring. The proof applies an
iterated theta operator to $\eta^{p+1}$ and uses Euler's pentagonal
identity to determine the sign. Combined with Robert's isogeny identity,
it gives
\[
\left(\frac{\Delta(\tau)}{\Delta(\ell\tau)}\right)^{(p^2-1)/24}
=\left(\frac{\ell}{p}\right)
\]
at every supersingular point of $X_0(\ell)$ for every prime $\ell\ne p$.
This determines the half-exponent sign in Ono's supersingular
discriminant-ratio identity and extends to arbitrary prime-to-$p$
isogenies. A companion paper applies the identity to the residues of
$T_N|U_p$ modulo $p$ for Fricke Hauptmoduln $T_N$.
\end{abstract}
\noindent\textbf{Keywords:} supersingular congruences; Eisenstein series; discriminant ratios; isogenies; theta operator.

\section{Introduction}\label{sec:intro}
Baker's supersingular congruence determines the square of a ratio
formed from $E_{p+1}$ and $\Delta$ \citep{Baker98}. We evaluate the
ratio itself explicitly, obtaining a factorial and a Legendre symbol
that are independent of the supersingular elliptic curve. This refines
Baker's congruence and, together with Robert's isogeny identity,
determines the half-exponent sign of a normalized discriminant ratio
along every prime-to-$p$ isogeny. Neither result requires a genus-zero
hypothesis.

A companion paper \citep{JordanFricke} uses Theorem~\ref{thm:C} below to
determine the signed residues of the reduction of $T_N|U_p$ modulo $p$,
where $T_N$ is the Hauptmodul of a genus-zero Fricke quotient
$X_0(N)^+$ of prime level.

\subsection{Notation}
Let $p\ge5$ be prime and put
\[
h=(p-1)/2,\qquad k=(p^2-1)/24,\qquad
A=\overline E_{p-1},\qquad B=\overline E_{p+1},\qquad
\Delta=\eta^{24}.
\]
Here $E_w$ is the normalized Eisenstein series of weight $w$, a bar
denotes reduction modulo $p$, and $A$ is the Hasse invariant. The
$q$-expansions of $A$ and $B$ are $1$ and $E_2$ modulo $p$,
respectively. We write $q=e^{2\pi i\tau}$, $\theta=q\,d/dq$,
$\Fp=\mathbb F_p$, $\Fq=\mathbb F_{p^2}$ and
$\Fb=\overline{\mathbb F}_p$.
The symbol $\left(\frac{a}{p}\right)$ denotes the Legendre symbol.
When a modular form is evaluated at an elliptic curve, a nonzero
invariant differential is understood. Ratios of equal weight are
independent of its choice. In particular, set
\[
c(E):=\frac{B(E,\omega)^h}{\Delta(E,\omega)^k}.
\]

\subsection{The factorial congruence and discriminant ratios}
The first result evaluates $c(E)$, rather than only its square.

\begin{theorem}[supersingular factorial evaluation]\label{thm:factorial}
For every prime $p\ge5$, every supersingular elliptic curve $E/\Fb$ and
every nonzero invariant differential $\omega$,
\begin{equation}\label{eq:factorial-evaluation}
B(E,\omega)^h=\Lp6\,h!\,\Delta(E,\omega)^k.
\end{equation}
Equivalently, in the graded ring $\Fp[E_4,E_6]$,
\begin{equation}\label{eq:factorial-divisibility}
A\mid B^h-\Lp6\,h!\,\Delta^k.
\end{equation}
\end{theorem}

The proof in Section~\ref{sec:factorial} applies the Serre--Robert
theta-iteration method to $F=\eta^{p+1}$. Euler's pentagonal identity
makes the relevant quadratic twist explicit, and comparison with the
iterated operator determines the factorial and its sign uniformly in $p$.
Squaring the evaluation recovers Baker's congruence.

\begin{corollary}[constancy and Baker's congruence]\label{cor:D}
The value $c(E)$ is the same for all supersingular $E/\Fb$, namely
$c_p=\Lp6\,h!\in\Fp^\times$. In particular
$c_p^2=-\Lp{-1}$, and
\[
E_{p+1}^{p-1}\equiv-\Lp{-1}\Delta^{(p^2-1)/12}
\pmod{(p,E_{p-1})}.
\]
\end{corollary}

For a prime $\ell$, let $Y_0(\ell)$ denote the non-cuspidal modular
curve over $\Z[1/\ell]$. For $x=(E,C)$, let $\pi:E\to E'=E/C$ and
choose a nonzero invariant differential $\omega'$ on $E'$. Define
\begin{equation}\label{eq:Dmoduli}
D_\ell(E,C):=\ell^{12}
\frac{\Delta(E,\pi^*\omega')}{\Delta(E',\omega')}.
\end{equation}
The differential weights make this independent of $\omega'$ and of the
representative of $x$. It is a unit on $Y_0(\ell)$ over $\Z[1/\ell]$,
with $q$-expansion $\Delta(q)/\Delta(q^\ell)$: on the Tate curve with
$C=\mu_\ell$, $\pi^*\omega'_{\rm can}=\ell\omega_{\rm can}$.
Its complex realization is therefore $\Delta(\tau)/\Delta(\ell\tau)$.
We use the same notation for its reduction at $p\ne\ell$.

\begin{samepage}
\begin{theorem}[signed discriminant ratio]\label{thm:C}
Let $p\ge5$ and $\ell\ne p$ be primes. At every supersingular point
$x$ of $Y_0(\ell)_{\Fb}$,
\[
D_\ell(x)^{(p^2-1)/24}=\leg{\ell}{p}.
\]
\end{theorem}
\end{samepage}



% TODO(author): theorem and remark numbers for Ono are those of arXiv v2
% (3 July 2026); numbering differs between versions, so recheck before submission.
\citet[Theorem~1.2]{Ono25} proves the corresponding identity at exponent
$2k=(p^2-1)/12$ for the reciprocal ratio
$\Delta(\ell\tau)/\Delta(\tau)$, and his Remark~4 exhibits values
with nontrivial half-exponent sign. Theorem~\ref{thm:C} determines
that sign uniformly as $\left(\frac{\ell}{p}\right)$.
Section~\ref{sec:proof} deduces the theorem from
Theorem~\ref{thm:factorial} and Robert's isogeny identity, and extends
it to arbitrary prime-to-$p$ isogenies. This is a refinement of the
supersingular identity used by Ono; no new assertion about the
Adams--Novikov spectral sequence is made here.

\subsection{Relation to earlier work}
\citet{Baker98} proves the supersingular congruence for $B^{p-1}$
recovered in Corollary~\ref{cor:D}. Theorem~\ref{thm:factorial}
refines it by evaluating the half-power ratio $c(E)$ explicitly and
showing that this value is independent of the supersingular curve.
The isogeny transformation of $B$ is due to
\citet[Th\'eor\`eme B, p.~106; \S3.3, Lemma~7 and (19), p.~124]{Robert80};
we include a $q$-expansion proof to fix the differential normalization.

The factorial evaluation uses a classical mechanism: Serre's functional
equation for Eisenstein series, presented by
\citet[\S2.5, Th\'eor\`eme~2, p.~119]{Robert80}, who attributes the relation
to a letter of Serre dated 13 June 1978 \citep[\S0.3, p.~109]{Robert80}, produces factorial
coefficients by iterating theta and then evaluating on the supersingular
locus. We apply that mechanism to the form $\eta^{p+1}$ of weight
$(p+1)/2$ at full level twelve. Its explicit quadratic twist gives the
signed value of $c_p$ simultaneously in all residue classes of $p$
modulo twelve. % Robert (1980) has been checked for this comparison: the citations to
% Theoreme B (p. 106), Theoreme 2 (p. 119), Lemme 7 and (19) (p. 124) and
% Theoreme 1 (p. 116) are correct, and the paper contains no evaluation of
% a power of E_{p+1} against a power of Delta. (Intro, Sec. 0.3, 2.3, 2.5,
% 3.3 and 5.4 were read; the remaining pages were only text-searched.)
% TODO(author): the same comparison is still unchecked against Serre's and
% Katz's work on E_{p+1} on the supersingular locus and Landweber (1988).
Robert's cited theorem concerns pairs of even-weight
Eisenstein series with complementary weights summing to $p+1$; it does
not state the factorial evaluation of Theorem~\ref{thm:factorial}. The argument here does not use
Baker's congruence or the signed discriminant-ratio identity as an input.

\citet[Theorem~3, p.~72; \S4]{Landweber88} gives the
half-exponent identity
$\Delta^{(p^2-1)/24}=(e_a-e_b)^{(p^2-1)/4}$ for supersingular
models $y^2=4\prod_{a=1}^3(x-e_a)$ and $a\ne b$.
Theorem~\ref{thm:C} instead compares discriminants along a prime-degree
isogeny, with its differential normalization explicit, and determines
the sign from Theorem~\ref{thm:factorial} and Robert's identity.
\citet[Theorem~1.2 and Remark~4]{Ono25}
proves the squared isogeny-ratio identity and records examples with nontrivial
half-exponent sign. His work develops the supersingular congruences arising in
the description of the divided beta family by \citet[Theorem~1.2]{Behrens09}
and the explicit calculations of \citet[\S5.4, Conjecture~5.17]{Larson20}.
These connections motivate the discriminant identity; no new statement about the
Adams--Novikov spectral sequence is needed here.

\paragraph{Organization.}
Section~\ref{sec:factorial} proves the factorial evaluation and
Corollary~\ref{cor:D}. Section~\ref{sec:proof} proves the signed
discriminant-ratio identity and its extension to arbitrary prime-to-$p$
isogenies. Section~\ref{sec:comp} records exact-arithmetic checks.

\paragraph{Acknowledgement of AI assistance.}
OpenAI GPT-6.1 Sol and Claude Opus 5.5 assisted in developing arguments, checking 
calculations and preparing this manuscript, including the factorial-evaluation 
argument and its exact-arithmetic verification code. The source computations
and proof checks are documented in the supplementary packages.
% TODO(author): confirm this statement and check the target journal's AI policy.
The author takes responsibility for the entire content of the paper.

\section{The supersingular factorial evaluation}\label{sec:factorial}
Put $h=(p-1)/2$ and $m=(p+1)/2$, so $m+h=p$ and
$m+h(p+1)=pm$. We prove Theorem~\ref{thm:factorial} by comparing two
expressions for an iterated theta operator applied to $F=\eta^{p+1}$.
The leading coefficient calculation is the one underlying Serre's
functional equation for Eisenstein series, as presented by
\citet[\S2.5, Th\'eor\`eme~2, p.~119]{Robert80}. The choice of $F$ makes
its quadratic twist explicit at every $p\ge5$.

\subsection{A form at full level twelve}
Work on a geometrically connected component $X(12)$ of the compactified
full-level modular curve, over $R=\Z_{(p)}[\zeta_{12}]$ localized at a
prime above $p$. This is a smooth proper curve carrying the Hodge line
bundle $\omega$. Forms of weight $w$ are sections of $\omega^{\otimes w}$;
a choice of nonzero invariant differential on an elliptic curve allows
evaluation of such a section. The cusp parameter is $q^{1/12}$.
We use the integral $q$-expansion principle on this prime-to-$p$ level
\citep[\S1.6]{Katz73}; see also \citealp{DR73}.

\begin{lemma}\label{lem:eta-level12}
The function $\eta^2$ is a holomorphic modular form of weight one with
trivial character on $\Gamma(12)$. In particular
$F=\eta^{p+1}\in M_m(\Gamma(12))$ is integral over $R$, and
\begin{equation}\label{eq:factorial-Fpowers}
F^{24}=\Delta^{p+1},\qquad F^{p-1}=\Delta^k.
\end{equation}
\end{lemma}
\begin{proof}
The eta transformation law \citep[Ch.~4]{Knopp70} gives
$\eta^2(\gamma\tau)=\nu(\gamma)^2(c\tau+d)\eta^2(\tau)$.
For $\gamma=\begin{pmatrix}a&b\\c&d\end{pmatrix}\in\mathrm{SL}_2(\Z)$ with
$c>0$ even, and hence $d$ odd, the multiplier is
\[
\nu(\gamma)=\leg{c}{|d|}
\exp\!\left(\frac{\pi i}{12}
\bigl((a+d)c-bd(c^2-1)+3d-3-3cd\bigr)\right),
\]
where the first factor is a Jacobi symbol. This case applies to every
$\gamma\in\Gamma(12)$ with $c>0$, since $12\mid c$. The symbol squares to
one. Since $a\equiv d\equiv1$ and $b\equiv c\equiv0\pmod{12}$, each of
$(a+d)c$, $bd(c^2-1)$, $3(d-1)$ and $3cd$ is divisible by $12$, so
\[
\nu(\gamma)^2=
\exp\!\left(\frac{2\pi i}{12}
\bigl((a+d)c-bd(c^2-1)+3d-3-3cd\bigr)\right)=1.
\] For $c<0$, apply the same formula to
$-\gamma$: its numerator is $-6$ modulo $12$, so its multiplier squared
is $-1$; the factor $-c\tau-d$ supplies the other minus sign.
For $c=0$, $\gamma$ is a translation by a multiple of $12$ and the
assertion follows from $\eta^2(\tau+12)=\eta^2(\tau)$.
The form is holomorphic at every cusp. The transformation law and the
product expansion give integral cusp expansions over $R$ (the multipliers
are twelfth roots of unity), so the $q$-expansion principle supplies the
integral section. Raising to the $m$th power gives $F$.
Both equalities in \eqref{eq:factorial-Fpowers} follow from
$\eta^{24}=\Delta$ and $24k=p^2-1$.
\end{proof}

\subsection{The iterated operator and the twist}
Let $\partial_w=\theta-\frac{w}{12}E_2$ be the Serre derivative.
By the quasi-modular transformation law of $E_2$, it sends modular forms
of weight $w$ on $\Gamma(12)$ to modular forms of weight $w+2$ on
$\Gamma(12)$. It also preserves $p$-integrality. Indeed, if $f$ has
expansion $\sum_na_nq^{n/12}$ with all $a_n\in R$, then
$\theta f=\sum_n(n/12)a_nq^{n/12}$ and $E_2f$ again have coefficients in
$R$, because $12$ is a unit in $R$; the $q$-expansion principle
\citep[\S1.6]{Katz73} then shows that $\partial_wf$ is defined over $R$.
Define
\[
G_0=F,\qquad
G_{i+1}=E_{p-1}\partial_{w_i}G_i+\frac{w_i}{12}E_{p+1}G_i,
\qquad w_i=m+i(p+1).
\]
Then $G_i$ is a $p$-integral form of weight $w_i$. Modulo $p$, one step of
this recursion is the operator $A\theta$ of \citet{Katz77}, which raises
the weight by $p+1$ and acts as $\theta$ on $q$-expansions. The Kummer congruences
$E_{p-1}\equiv1$ and $E_{p+1}\equiv E_2\pmod p$ show that its
$q$-expansion is $\theta^iF\pmod p$. Moreover
\begin{equation}\label{eq:factorial-leading}
G_i=\gamma_i E_{p+1}^{\,i}F+E_{p-1}R_i,\qquad
\gamma_i=12^{-i}\prod_{j=0}^{i-1}w_j,
\end{equation}
with $R_i$ a $p$-integral modular form. This follows by induction: in the
recursion the derivative term is multiplied by $E_{p-1}$, and the
remaining term multiplies the leading coefficient by $w_i/12$.
At $i=h$, Wilson's theorem gives
\begin{equation}\label{eq:factorial-gamma}
\overline\gamma_h
=12^{-h}\prod_{j=0}^{h-1}(m+j)
=12^{-h}\frac{(p-1)!}{h!}
=-\frac{12^{-h}}{h!}\quad\text{in }\Fp.
\end{equation}

Write $\chi_{12}(n)=\leg{12}{n}$ for the Kronecker character, with
values $1,-1,0$ on $n\equiv\pm1,\pm5$, and nonunits modulo $12$,
respectively. Euler's pentagonal identity \citep{Apostol90} takes the form
\[
\eta(\tau)=\sum_{n\ge1}\chi_{12}(n)q^{n^2/24}.
\]
\begin{lemma}[quadratic twist]\label{lem:factorial-twist}
The following identity holds between mod-$p$ forms of weight $pm$ on
$\Gamma(12)$:
\begin{equation}\label{eq:factorial-twist}
\overline G_h=\Lp6\bigl(A^m\overline F-\chi_{12}(p)\overline F^{\,p}\bigr).
\end{equation}
\end{lemma}
\begin{proof}
By the product expansion and the freshman's dream,
$F\equiv\eta(\tau)\eta(p\tau)\pmod p$. Thus
\[
\overline F=\sum_{a,b\ge1}\chi_{12}(a)\chi_{12}(b)
q^{(a^2+pb^2)/24}.
\]
Applying $\theta^h$ multiplies each term by
$((a^2+pb^2)/24)^h$, which is zero modulo $p$ when $p\mid a$
and equals $\leg{24}{p}=\Lp6$ otherwise.
The sub-sum with $a=pa'$ is
\[
\chi_{12}(p)\sum_{a',b\ge1}\chi_{12}(a')\chi_{12}(b)
q^{p(b^2+pa'^2)/24}
=\chi_{12}(p)\overline F^{\,p}.
\]
Consequently the $q$-expansion of $\overline G_h$ equals
$\Lp6(\overline F-\chi_{12}(p)\overline F^{\,p})$.
The integral forms $G_h$, $E_{p-1}^mF$ and $F^p$ all have weight
$pm$. Since $A$ has $q$-expansion one, the $q$-expansion principle
proves \eqref{eq:factorial-twist} as an identity of sections.
\end{proof}

\begin{proof}[Proof of Theorem~\ref{thm:factorial}]
Let $(E,\omega_E)$ be supersingular and choose a full level-$12$
structure $\alpha$ with the fixed Weil pairing. Evaluate
\eqref{eq:factorial-leading} and \eqref{eq:factorial-twist} at
$x=(E,\alpha,\omega_E)$. Since $A(x)=0$, they give
\[
-\frac{12^{-h}}{h!}\,B(E,\omega_E)^h\overline F(x)
=-\Lp6\chi_{12}(p)\overline F(x)^p.
\]
Equation~\eqref{eq:factorial-Fpowers} shows that
$\overline F(x)\ne0$ and
$\overline F(x)^{p-1}=\Delta(E,\omega_E)^k$. Hence
\[
B(E,\omega_E)^h
=\Lp6\chi_{12}(p)12^h h!\,\Delta(E,\omega_E)^k.
\]
Euler's criterion gives $12^h=\leg{12}{p}=\chi_{12}(p)$,
so the two character factors cancel. This proves the pointwise formula
for every supersingular curve, including non-rational $j$-invariants.

For the divisibility statement, work in $\Fb[E_4,E_6]$.
The Hasse invariant $A$ is squarefree and prime to $B$
\citep[\S2.3, Th\'eor\`eme~1(ii), p.~116]{Robert80}. Its factors are $E_4$ if
$j=0$ is supersingular, $E_6$ if $j=1728$ is supersingular, and
$E_4^3-j_0\Delta$ for the remaining supersingular $j_0$
\citep[\S3]{KZ98}. This is the geometric factorization over $\Fb$;
the irreducible quadratic factors in $j$ over $\Fp$ in
\citet[Theorem~1.2]{Baker98} split after extending the constants.
For $j_0\ne0,1728$, the factor
\[
E_4^3-j_0\Delta
=\left(1-\frac{j_0}{1728}\right)E_4^3
+\frac{j_0}{1728}E_6^2
\]
is irreducible: it has the form $aX^3+bY^2$ with $ab\ne0$.
The zero set of each factor away from the origin
parametrizes the corresponding curves with their nonzero differentials.
The pointwise identity therefore forces each factor to divide
$B^h-\Lp6h!\Delta^k$, so their product $A$ divides it.
Divisibility descends to $\Fp[E_4,E_6]$. Conversely, divisibility
immediately implies the pointwise identity because $A$ vanishes at every
supersingular curve. The factors $E_4$ and $E_6$ account for the elliptic
points $j=0,1728$ without extra multiplicities in the graded ring.
\end{proof}

\begin{proof}[Proof of Corollary~\ref{cor:D}]
Theorem~\ref{thm:factorial} gives $c(E)=\Lp6h!\in\Fp^\times$
independently of $E$. Squaring and using
$(h!)^2=(-1)^{h+1}=-\Lp{-1}$ proves the stated form of
Baker's congruence. This proof uses neither that congruence nor
connectivity of supersingular isogeny graphs as an input.
\end{proof}

\begin{remark}[the constant $c_p$]\label{rem:mordell}
The constant $c_p=\Lp6\,h!$ is not an elementary sign. If
$p\equiv1\pmod4$, then $c_p^2=-1$, so $c_p$ is a square root of $-1$ in
$\Fp$. If $p\equiv3\pmod4$, then $c_p=\pm1$, and the sign is governed by a
class number: writing $\mathrm{cl}(-p)$ for the class number of
$\Q(\sqrt{-p})$, Mordell's congruence
$h!\equiv(-1)^{(\mathrm{cl}(-p)+1)/2}\pmod p$ \citep{Mordell61} gives
\[
c_p=\Lp6\,(-1)^{(\mathrm{cl}(-p)+1)/2}\qquad(p\equiv3\pmod4).
\]
For example, $c_7=c_{11}=c_{23}=1$ and $c_{19}=c_{31}=-1$.
\end{remark}

\begin{remark}[a limit to higher-root constancy]\label{rem:higher-roots}
Constancy in Corollary~\ref{cor:D} need not persist for a further
square root. For $p\equiv1\pmod8$, the ratio
$\rho_p=B^{(p-1)/4}/\Delta^{(p^2-1)/48}$ has weight zero and satisfies
$\rho_p^2=c_p$ at supersingular points, but it need not be constant.
For example, modulo $17$,
\[
A=8E_4^4+10E_4E_6^2,\qquad B=7E_4^3E_6+11E_6^3.
\]
The two pairs of form values below have $A=0$ and $\Delta\ne0$, and
hence represent supersingular elliptic curves with nonzero
invariant differentials. Using $\Delta=(E_4^3-E_6^2)/1728$ gives
\[
\begin{array}{c|cc|cc|c}
j&E_4&E_6&B&\Delta&\rho_{17}=B^4/\Delta^6\\\hline
0&0&1&11&3&-2\\
8&3&3&14&14&2
\end{array}
\qquad\text{in }\mathbb F_{17}.
\]
Both values square to $c_{17}=4$. Thus a uniformly constant
quarter-power analogue is false. More generally, whenever
$r\mid h$ and $r\mid k$, the ratio $B^{h/r}/\Delta^{k/r}$ takes at
most $r$ values, since its $r$th power is $c_p$; this bound alone
does not imply that all $r$ values occur.
\end{remark}

\section{Signed discriminant ratios}\label{sec:proof}
The factorial evaluation makes $c(E)$ constant on the supersingular
locus. Robert's isogeny identity compares $B$ on the two curves of an
isogeny and turns this constancy into a formula for their discriminants.
We first record the differential normalization.
\begin{lemma}[Robert's isogeny identity]\label{lem:Bisog}
Let $\ell $ be any prime $\ne p$. Let $x=(E,C)$ be a supersingular point of
$X_0(\ell )$, $\pi:E\to E'=E/C$, and $\omega=\pi^*\omega'$. Then
$B(E',\omega')=\ell \,B(E,\omega)$, $A(E',\omega')=A(E,\omega)$ and
$\Delta(\ell \tau)(x)=\ell ^{-12}\Delta(E',\omega')$.
\end{lemma}
\begin{proof}
The isogeny identity is \citet[Th\'eor\`eme B; \S3.3, Lemma~7]{Robert80}; see also
\citet[Lemma~2.2]{Baker98}. We verify the normalization by $q$-expansions.
For a level-one form $f$ of weight $w$, $f(\ell \tau)$ at $(E,C,\omega)$ equals $\ell ^{-w}f(E',\omega')$. Check
this on $(\mathrm{Tate}(q),\mu_\ell )$, where $E'\cong\mathrm{Tate}(q^\ell )$ and
$\pi^*\omega_{\rm can}'=\ell \omega_{\rm can}$. The form $E_2(\tau)-\ell E_2(\ell \tau)$ is a holomorphic
weight-$2$ form on $\Gamma_0(\ell )$ with $p$-integral coefficients. Modulo $p$, the forms
$B(\tau)-\ell B(\ell \tau)$ and $A\cdot(E_2-\ell E_2(\ell \tau))$ have weight $p+1$ and the same $q$-expansion, so they
are equal. At a supersingular point this gives $B(E,\omega)=\ell \cdot \ell ^{-(p+1)}B(E',\omega')=\ell ^{-1}B(E',\omega')$,
because $\ell ^{p-1}=1$. The statement about $A$ follows in the same way from $A(\ell \tau)=A(\tau)$.
\end{proof}

We use the algebraic unit $D_\ell$ of \eqref{eq:Dmoduli}, with
$\pi:E\to E'=E/C$ and $\omega=\pi^*\omega'$.

\begin{proof}[Proof of Theorem~\ref{thm:C}]
At a supersingular point $x=(E,C)$, Lemma~\ref{lem:Bisog} and
\eqref{eq:Dmoduli} imply
\[
\frac{c(E')}{c(E)}
=\ell^h\left(\frac{\Delta(E,\omega)}{\Delta(E',\omega')}\right)^k
=\Lp\ell\,D_\ell(x)^k,
\]
because $\ell^{12k}=\ell^{(p^2-1)/2}=1$.
By Corollary~\ref{cor:D}, the left side is one and is nonzero.
Therefore $D_\ell(x)^k=\Lp\ell$.
\end{proof}

\begin{corollary}[arbitrary prime-to-$p$ isogenies]\label{cor:degree-n}
Let $\varphi:E\to E'$ be an isogeny of degree $n\ge1$ between
supersingular elliptic curves over $\Fb$, with $(n,p)=1$. For any
nonzero invariant differential $\omega'$ on $E'$, put
\[
D_\varphi:=n^{12}
\frac{\Delta(E,\varphi^*\omega')}{\Delta(E',\omega')}.
\]
Then $D_\varphi$ is independent of $\omega'$ and
\[
D_\varphi^k=\Lp n.
\]
In particular, for cyclic kernels the signed discriminant-ratio
identity holds on $Y_0(n)_{\Fb}$ for every $n$ prime to $p$.
\end{corollary}
\begin{proof}
Over $\Fb$, factor $\varphi$ into isogenies of prime degree and apply
Lemma~\ref{lem:Bisog} successively, choosing differentials by pullback.
This gives $B(E',\omega')=nB(E,\varphi^*\omega')$. Consequently
\[
1=\frac{c(E')}{c(E)}
=n^h\left(\frac{\Delta(E,\varphi^*\omega')}{\Delta(E',\omega')}\right)^k
=\Lp n\,D_\varphi^k,
\]
since $n^{12k}=1$ in $\Fp$. The Legendre symbol is its own inverse.
The weight of $\Delta$ gives independence of the chosen differential.
\end{proof}

\begin{remark}[eta coordinates]\label{rem:eta}
For $N\in\{2,3,5,7,13\}$ put $a_N=24/(N-1)$ and
$u_N=(\eta(\tau)/\eta(N\tau))^{a_N}$, a coordinate on $X_0(N)$. Then
$D_N=u_N^{N-1}$, and Theorem~\ref{thm:C} gives
\[
u_N(x)^{(p^2-1)/a_N}=\Lp N
\]
at every supersingular point $x$ of $Y_0(N)_{\Fb}$, for every prime
$p\ge5$ with $p\ne N$. For example, $j=(u_2+256)^3/u_2^2$ on $X_0(2)$.
For $p=5$ the only supersingular invariant is $j=0$, attained at
$u_2=-256\equiv-1=\leg25$. For $p=7$ it is $j=1728\equiv6$, attained
exactly at $u_2\equiv\pm1$, and $u_2^2=1=\leg27$. This is the form in
which the identity is used in the companion paper \citep{JordanFricke}.
\end{remark}

\begin{remark}[comparison with Ono]\label{rem:ono}
\citet[Theorem~1.2 and Remark~4]{Ono25} uses the reciprocal
$R_\ell=\Delta(\ell\tau)/\Delta(\tau)=D_\ell^{-1}$.
Since a Legendre symbol equals its inverse, Theorem~\ref{thm:C}
determines the same half-exponent sign for $R_\ell$, namely $\Lp\ell$.
For example, $\leg2{11}=-1$, as in his example.
His Theorem~1.3 states that on a model with $p^2$-Frobenius $[-p]$
and an $\Fq$-rational differential,
$\Delta(E,\omega)^{2k}=(-1)^{(p+1)/2}$.
Theorem~\ref{thm:factorial} instead relates $B(E,\omega)^h$ to
$\Delta(E,\omega)^k$ for every nonzero differential.
The ratio argument requires no descent of a twenty-fourth root of
$\Delta$ to $\Fq$.
\end{remark}

\section{Computational evidence}\label{sec:comp}
The following exact-arithmetic computations are checks on the theorems,
not inputs to their proofs.

For Theorem~\ref{thm:factorial}, exact polynomial remainders and evaluations
at all 6430 supersingular $j$-values for the 166 primes $5\le p\le997$
(4734 outside $\Fp$) agree with the asserted constant. Evaluations use three
randomly rescaled models per value, with additional Legendre and Weierstrass
coefficient routes.
The two theta identities in the proof were also checked in
$\Fp[E_4,E_6]$ at $p=23,47,71,167$, where $\eta^{p+1}$ has level one.

Theorem~\ref{thm:C} was tested independently using normalized
V\'elu quotients and discriminants, without eta functions or Hauptmoduln.
For the 15 primes $\ell\le47$ and admissible $5\le p\le199$, all 647 pairs
and 127541 cyclic-kernel checks gave the asserted discriminant-ratio sign.
These count kernels on chosen supersingular representatives, not distinct
coarse points, since automorphisms can identify kernels at $j=0,1728$.

\paragraph{Data and software.}
The supplementary material for this paper consists of three packages.
Their folder names are those of the archived release; they do not reflect
the current status or numbering of the results.
\begin{itemize}[leftmargin=*,itemsep=1pt]
\item \path{Factorial_Conjecture_Result_v1_0}: code and recorded results
for Theorem~\ref{thm:factorial} under \texttt{new\_work/}, with
reproduction commands in \texttt{run\_all.sh} (\texttt{cypari2} 2.2.4,
PARI/GP 2.17.2 and \texttt{mpmath} 1.3.0).
\item \path{Eta_Frobenius_Lemma_Audit_v1_0}: the V\'elu-quotient checks
of Theorem~\ref{thm:C}.
\item \path{Theorem_1_3_Benchmark_v1_0}: timings for exact evaluation of
$B^h$ by Theorem~\ref{thm:factorial} (numbered Theorem~1.3 in an earlier
draft), with the timing code, exact inputs and outputs, raw timing
samples and verification records.
\end{itemize}
\if\relax\detokenize\expandafter{\DataRepositoryURL}\relax
The code, per-pair data and reproduction instructions are supplied as
supplementary files with the manuscript.
\else
The code, data and reproduction instructions are publicly available at
\expandafter\url\expandafter{\DataRepositoryURL}.
\fi

\section{Further directions}
The proof of Theorem~\ref{thm:factorial} suggests seeking other families
of modular forms whose quadratic twists admit similarly explicit
expressions. Extending this proof strategy would require a suitable
analogue of the twist identity, together with the necessary weight,
integrality and nonvanishing properties.
Remark~\ref{rem:higher-roots} shows that further root extraction need
not preserve constancy: the half-power evaluation cannot in general
be replaced by a uniformly constant quarter-power analogue.

The discriminant-ratio identity holds without a genus-zero hypothesis,
including for arbitrary prime-to-$p$ isogenies. Its application to
Fricke Hauptmoduln, which additionally requires an integral genus-zero
coordinate, is the subject of the companion paper \citep{JordanFricke}.

\section*{Code and data availability}
All code, scripts and logs are available in the GitHub repository
\begin{center}\url{https://github.com/bbrhuft/factorial-supersingular-congruence}\end{center}
The computational materials described here are those of release v2.0 of this repository.
Earlier versions of this work, in which this paper and its companion formed a single
manuscript, are archived on Zenodo at
\href{https://doi.org/10.5281/zenodo.24024359}{doi:10.5281/zenodo.24024359} (v1.11) and
\href{https://doi.org/10.5281/zenodo.24084789}{doi:10.5281/zenodo.24084789} (v1.13).

\begin{thebibliography}{99}\small
\bibitem[Apostol(1990)]{Apostol90} Apostol, T.~M. (1990). \emph{Modular Functions and Dirichlet Series in Number Theory}, 2nd ed. Graduate Texts in Mathematics 41. Springer.


\bibitem[Baker(1998)]{Baker98} Baker, A. (1998). A supersingular congruence for modular forms.
\emph{Acta Arithmetica}, 86(1), 91--100. \url{https://doi.org/10.4064/aa-86-1-91-100}.


\bibitem[Behrens(2009)]{Behrens09} Behrens, M. (2009). Congruences between modular forms given by the divided $\beta$ family in homotopy theory. \emph{Geometry \& Topology}, 13, 319--357. \url{https://doi.org/10.2140/gt.2009.13.319}.


\bibitem[Deligne and Rapoport(1973)]{DR73} Deligne, P., \& Rapoport, M. (1973). Les sch\'emas de modules de
courbes elliptiques. In \emph{Modular Functions of One Variable II} (LNM 349, pp.~143--316). Springer.


\bibitem[Jordan(2026)]{JordanFricke} Jordan, D. (2026). Supersingular poles and signed residues of
prime-level Fricke Hauptmoduln. Preprint (companion paper). Manuscript, code and data: \url{https://github.com/bbrhuft/fricke-signed-residues}.

\bibitem[Kaneko and Zagier(1998)]{KZ98} Kaneko, M., \& Zagier, D. (1998). Supersingular $j$-invariants,
hypergeometric series, and Atkin's orthogonal polynomials. In \emph{Computational Perspectives on Number
Theory} (AMS/IP Studies in Advanced Mathematics 7, pp.~97--126).


\bibitem[Katz(1973)]{Katz73} Katz, N.~M. (1973). $p$-adic properties of modular schemes and modular forms.
In \emph{Modular Functions of One Variable III} (LNM 350, pp.~69--190). Springer.


\bibitem[Katz(1977)]{Katz77} Katz, N.~M. (1977). A result on modular forms in characteristic $p$.
In \emph{Modular Functions of One Variable V} (LNM 601, pp.~53--61). Springer.


\bibitem[Knopp(1970)]{Knopp70} Knopp, M.~I. (1970). \emph{Modular Functions in Analytic Number Theory}. Markham.


\bibitem[Landweber(1988)]{Landweber88} Landweber, P.~S. (1988). Supersingular elliptic curves and congruences for Legendre polynomials. In \emph{Elliptic Curves and Modular Forms in Algebraic Topology}, Lecture Notes in Mathematics 1326, pp.~69--93. Springer.


\bibitem[Larson(2020)]{Larson20} Larson, D.~M. (2020). Explicit modular forms from the divided beta family. \emph{Topology and its Applications}, 273, 107094. \url{https://doi.org/10.1016/j.topol.2020.107094}.


\bibitem[Mordell(1961)]{Mordell61} Mordell, L.~J. (1961). The congruence $((p-1)/2)!\equiv\pm1\pmod p$.
\emph{American Mathematical Monthly}, 68, 145--146.


\bibitem[Ono(2026)]{Ono25} Ono, K. (2026). Modular forms for chromatic homotopy: supersingular congruences. \emph{arXiv:2509.16175v2}, 3 July 2026. \url{https://arxiv.org/abs/2509.16175v2}.


\bibitem[Robert(1980)]{Robert80} Robert, G. (1980). Congruences entre s\'eries d'Eisenstein, dans le cas supersingulier. \emph{Inventiones Mathematicae}, 61, 103--158.

\end{thebibliography}
	\bigskip
\noindent{\scshape Maynooth University Social Sciences Institute (MUSSI), Maynooth University, Maynooth, County Kildare, Ireland}\\
\emph{Email address}: \texttt{David.Jordan@mu.ie}
\end{document}

